%% Dokumentenklasse (KOMA-Script) ==========================================
%% scrreprt: Report-Klasse mit Chaptern, ideal für Abschlussarbeiten
\documentclass[%
   final,%              % Finales Dokument (vs. draft)
   paper=a4,%           % DIN A4 Papierformat
   pagesize=auto,%      % PDF-Seitengröße automatisch einstellen
   fontsize=12pt,%      % Schriftgröße 12pt (Standard für Abschlussarbeiten)
   version=last,%       % Neueste KOMA-Script Version
   parskip=half%        % Absatzabstand statt Einzug (half = halbe Zeile)
]{scrreprt}


%% Spracheinstellung ========================================================
%% Options: ngerman (Deutsch, neue Rechtschreibung), english (Englisch)
\def\lang{ngerman}

%% KI-Tools Declaration (erforderlich bei KI-Nutzung) =======================
%% Dokumentieren Sie verwendete KI-Tools hier und im Anhang "Verzeichnis der KI-Nutzung"
\def\declareUseOfGenerativeAITool{GitHub Copilot, ChatGPT}

%% Preambel laden ===========================================================
\input{preambel/settings}
\input{preambel/preambel}

% URLs im normalen Textfont statt Monospace anzeigen
\urlstyle{same}

% Setzt das Zugriffsdatum (urldate) automatisch auf das aktuelle Kompilierdatum (ISO-Format)
\DeclareSourcemap{
  \maps[datatype=bibtex]{
    \map[overwrite]{
      \pertype{online}
      \pertype{misc}
      \pertype{software}
      \step[fieldset=urldate, fieldvalue=\the\year-\ifnum\month<10 0\fi\the\month-\ifnum\day<10 0\fi\the\day]
    }
    % Entfernt redundante "Accessed:"-Einträge aus dem note-Feld automatisch
    \map{
      \step[fieldsource=note, match=\regexp{Accessed:}, final]
      \step[fieldset=note, null]
    }
  }
}

% "Besucht am" im Literaturverzeichnis immer kleinschreiben (erzwingt Kleinschreibung auch nach Punkt)
\DeclareFieldFormat{urldate}{\mkbibparens{besucht am\space#1}}

%% Metadaten der Arbeit ====================================================
%% Bitte alle Felder ausfüllen für korrekte Anzeige im PDF
\author{Maria Musterfrau}
\studentID{1234567}
\studentAddress{Musterstraße 1, 12345 Musterstadt}
\thesis{Bachelor-Thesis}
\title{Entwicklung und Evaluation agentenbasierter Systeme im Software Engineering: Ein Ansatz zur Automatisierung von Entwicklungsprozessen mittels Large Language Models}
\academicTitle{Bachelor of Science}
\firstReferee{Prof.~Dr.~Klaus Quibeldey-Cirkel}
\secondReferee{Prof.~Dr.~Philipp Diebold}

%%% Silbentrennung
\input{preambel/Hyphenation}

%% Dokument Beginn %%%%%%%%%%%%%%%%%%%%%%%%%%%%%%%%%%%%%%%%%%%%%%%%%%%%%%%%

\begin{document}
% Temporär pageanchor deaktivieren für Titelseite und Abstract (verhindert duplicate identifier)
\hypersetup{pageanchor=false}

% Deckblatt
\input{content/00_Titel}

% Frontmatter beginnt hier (römische Seitenzahlen)
\frontmatter{}

% Abstract - Kurzfassung der Arbeit ohne Wertung
% !TEX root = ../Thesis.tex

%% Leitfragen für das Abstract:
%% - Motivation: Warum ist das Thema relevant?
%% - Fragestellung: Was untersucht die Arbeit?
%% - Methode: Wie wurde untersucht (Ansatz, Techniken)?
%% - Ergebnisse: Zu welchen Schlussfolgerungen kam die Forschung?
%% - Relevanz: Was trägt die Arbeit zum bestehenden Wissen bei?
%%
%% Format: 150-250 Wörter, prägnant und eigenständig verständlich

\chapter*{Kurzfassung}
\addcontentsline{toc}{chapter}{Kurzfassung}

% Reduziere den Absatzabstand lokal, damit alles auf eine Seite passt
\begingroup
\setlength{\parskip}{0.5em}
\textbf{Kontext und Motivation:} Large Language Models (LLMs) entwickeln sich zunehmend von reinen Textgeneratoren zu agentischen Systemen, die planen, Werkzeuge einbinden und mehrschrittige Aufgaben bearbeiten können. Für das Software Engineering eröffnet dies neue Möglichkeiten zur Unterstützung von Code-Review, Testgenerierung, Fehlersuche und Refactoring.

\textbf{Zielsetzung:} Die Arbeit untersucht die Entwicklung und Evaluation agentischer Architekturen für Software-Engineering-Aufgaben. Im Zentrum stehen drei Fragen: (1) Wie lassen sich robuste Agentenstrategien für SE-Workflows systematisch modellieren? (2) Welche Architekturprinzipien ermöglichen eine sichere und nachvollziehbare Werkzeugintegration? (3) Mit welchen Metriken kann die Leistungsfähigkeit agentischer Systeme unter realistischen Randbedingungen bewertet werden?

\textbf{Methodik:} Auf Basis einer strukturierten Literaturanalyse wird eine Referenzarchitektur entwickelt, die Planung, Werkzeugnutzung, episodisches Gedächtnis und Sicherheitsmechanismen integriert. Darauf aufbauend wird ein prototypisches System in Python implementiert und anhand reproduzierbarer Szenarien aus dem Software Engineering evaluiert.

\textbf{Ergebnisse:} Die Evaluation zeigt für das gewählte Szenarienset, dass die entwickelte Architektur bei automatisiertem Refactoring eine Erfolgsrate von \pct{73} erreicht und die mittlere Bearbeitungszeit um \pct{45} reduziert. Reflexionsmechanismen verbessern die Robustheit gegenüber fehlerhaften Werkzeugausgaben um \pct{28}; optimierte Kontextverwaltung senkt die Token-Nutzung um bis zu \pct{40}. Die Resultate sind als indikative Ergebnisse einer prototypischen Umsetzung zu interpretieren und nicht als unmittelbar auf öffentliche Leaderboards übertragbare Vergleichswerte.

\textbf{Beitrag:} Die Arbeit liefert eine praxisnahe Referenzarchitektur, konkrete Implementierungsrichtlinien und ein mehrdimensionales Evaluationsschema für agentische SE-Systeme. Sie zeigt technische Potenziale, benennt methodische Grenzen und ordnet die Ergebnisse in die aktuelle Forschung zu agentischer Softwareentwicklung ein.%
\par
\endgroup

% Englisches Abstract (Standard in der Informatik)
\chapter*{Abstract}
\addcontentsline{toc}{chapter}{Abstract}

\begingroup
\setlength{\parskip}{0.5em}
\textbf{Context and Motivation:} Large Language Models (LLMs) are increasingly evolving from pure text generators into agentic systems that can plan, invoke tools, and handle multi-step tasks. In software engineering, this creates new opportunities to support code review, test generation, debugging, and refactoring.

\textbf{Objective:} This thesis investigates the development and evaluation of agentic architectures for software engineering tasks. The central questions are: (1)~How can robust agent strategies for SE workflows be modeled systematically? (2)~Which architectural principles enable secure and traceable tool integration? (3)~Which metrics allow a realistic assessment of agentic systems under practical constraints?

\textbf{Methodology:} Based on a structured literature review, a reference architecture is developed that integrates planning, tool usage, episodic memory, and safety mechanisms. A prototypical system is implemented in Python and evaluated using reproducible software-engineering scenarios.

\textbf{Results:} For the selected evaluation scenarios, the proposed architecture achieves a success rate of \pct{73} in automated refactoring and reduces mean processing time by \pct{45}. Reflection mechanisms improve robustness against faulty tool outputs by \pct{28}, while optimized context management reduces token usage by up to \pct{40}. These values should be interpreted as indicative prototype results rather than as directly comparable public leaderboard scores.

\textbf{Contribution:} The thesis provides a practical reference architecture, concrete implementation guidelines, and a multidimensional evaluation scheme for agentic SE systems. It identifies technical opportunities, names methodological limitations, and situates the findings within the current research landscape on agentic software engineering.%
\par
\endgroup

\cleardoublepage

\cleardoublepage{}

% Danksagung
\input{content/00_Danksagung}

% Hinweis zur geschlechtergerechten Sprache
\input{content/00_GenderErklaerung}

% Ab hier pageanchor aktivieren für korrekte PDF-Links
\hypersetup{pageanchor=true}

\cleardoublepage{}
% Inhaltsverzeichnis in den PDF-Links eintragen
\pdfbookmark[1]{Inhaltsverzeichnis}{toc}
\tableofcontents
\cleardoublepage%

% Abbildungs- und Tabellenverzeichnis
\listoffigures%
\listoftables%
\lstlistoflistings
\cleardoublepage

% Abkürzungsverzeichnis
\input{content/00_Abkuerzungen}
\cleardoublepage%

% Hauptteil beginnt hier (arabische Seitenzahlen ab 1)
\mainmatter{}
\chapter{Einführung}
\index{Agentic AI}%
\index{Software Engineering}%

% Längeres abgesetztes Zitat
\begin{quote}
\textit{\textquote[\cite{karpathy2023intro}]{Wir erleben die Entstehung eines neuen Betriebssystems: Das LLM fungiert als Kernel-Prozess, das Kontextfenster als Arbeitsspeicher und externe Werkzeuge als E/A-Schnittstellen. Dies markiert einen fundamentalen Wandel in der Art und Weise, wie wir Computerbefehle komponieren und Softwarearchitekturen denken.}}
\smallskip
\begin{flushright}
— Andrej Karpathy, \textit{Intro to Large Language Models}, 2023
\end{flushright}
\end{quote}

\section{Motivation und Relevanz}

Im Software Engineering zeichnet sich derzeit eine Verschiebung der Automatisierungslogik ab: \emph{agentische KI} \textendash{} also KI-Systeme, die Ziele interpretieren, Pläne erstellen, Werkzeuge verwenden und Ergebnisse eigenständig prüfen \textendash{} erweitert klassische Automatisierung um adaptive, mehrschrittige Problemlösung.\footnote{Der Begriff \emph{agentische KI} bezeichnet Systeme, die über mehrere Schritte hinweg komplexe Aufgaben bearbeiten und dabei Planung, Werkzeugnutzung und Rückkopplung kombinieren.}
 
Darüber hinaus verlangt die praktische Anwendung agentischer Systeme einen belastbaren Ausgleich zwischen Automatisierung und menschlicher Kontrolle. In vielen Entwicklungsprozessen sind Schritte, die bislang überwiegend manuell erfolgen \textendash{} etwa Code-Reviews, Testgenerierung oder Release-Checks \textendash{} geeignete Kandidaten für assistierte Abläufe. Agenten können Routinearbeit übernehmen, erzeugen damit jedoch zugleich neue Anforderungen an Validierung, Protokollierung und Eingriffsmöglichkeiten. Die vorliegende Arbeit untersucht diesen Zielkonflikt und entwickelt Gestaltungsprinzipien für einen kontrollierten produktiven Einsatz.
Wie Richards und Ford betonen, \textquote[\cite{richards2020fundamentals}]{müssen Architekturprinzipien auf Skalierbarkeit und Wartbarkeit ausgelegt sein}, um praktischen Anforderungen gerecht zu werden.
Für Forschung und Praxis ergeben sich daraus neue Architektur- und Methodikfragen: Wie lassen sich robuste agentische Workflows entwerfen? Wie werden Werkzeugnutzung, Gedächtnis und Langkontext so orchestriert, dass nachvollziehbare Entscheidungen entstehen? Und wie lassen sich Sicherheit, Prüfbarkeit und Testbarkeit systematisch in agentische Systeme integrieren?\footnote{Praktische Orchestrierung bezeichnet hier die Koordination von Planungsschritten, Werkzeugaufrufen und Gedächtniszugriffen in einer strukturierten Abfolge.}

\section{Problemstellung}

Vor dem Hintergrund adressiert die Arbeit exemplarisch die Entwicklung und Evaluation eines agentischen Systems für Softwareentwicklungsaufgaben (z.\,B. Refactoring, Code-Review, Generierung von Tests). Zentrale Fragen sind:

\begin{itemize}
  \item Wie lassen sich Agentenstrategien (Planen, Werkzeugaufrufe, Selbstkritik) systematisch modellieren?
  \item Wie werden externe Werkzeuge (VCS, CI, Linter, Issue-Tracker) sicher und nachvollziehbar eingebunden?
  \item Welche Metriken messen Fortschritt, Qualität und Sicherheit realistisch?
\end{itemize}
 
Die hier formulierten Probleme sind sowohl konzeptioneller als auch praktischer Natur: Im Mittelpunkt stehen nicht nur die abstrakte Modellierung von Strategien, sondern ebenso konkrete Implementierungsfragen (Schnittstellen, Serialisierung, Fehlerbehandlung) und Fragen des Evaluationsdesigns (Benchmarks, Messgrößen, Reproduzierbarkeit). Die Ergebnisse sollen Entwicklungsteams konkrete Handlungsempfehlungen geben, wie agentische Automatisierung schrittweise, kontrolliert und messtechnisch abgesichert eingeführt werden kann.\index{Evaluation}
\section{Zielsetzung}

Die Ziele der Arbeit sind auf die Spezialisierung \enquote{Software Engineering mit agentischer KI} zugeschnitten:\footnote{Siehe auch~\cite{wang2023survey} für verwandte Arbeiten zu Agentenarchitekturen.}

\begin{enumerate}
  \item Analyse des Forschungsstands zu agentischen Architekturen und Orchestrierungsframeworks
  \item Entwurf einer referenzierbaren Agentenarchitektur für Software-Engineering-Aufgaben
  \item Implementierung eines prototypischen Agenten mit Werkzeuganbindung und Gedächtnis
  \item Evaluation anhand reproduzierbarer Benchmarks (Qualität, Kosten, Laufzeit, Sicherheit)
\end{enumerate}
 
Kurz gefasst zielt die Arbeit darauf ab, aus Praxisproblemen generalisierbare Lösungen abzuleiten: Neben einem lauffähigen Prototyp steht die Frage im Vordergrund, welche Architekturmuster sich zuverlässig übertragen lassen und welche Messgrößen den Mehrwert gegenüber manuellen Prozessen belastbar quantifizieren. Damit richtet sich die Arbeit gleichermaßen an Forschende und an Verantwortliche in der Softwareentwicklung.
\section{Abgrenzung des Themas}

Die Arbeit fokussiert Agenten für Softwareentwicklungsaufgaben. Nicht im Fokus sind z.\,B. Reinforcement Learning from Human Feedback (RLHF) im Detail, Trainingsmethoden auf Rohdaten oder proprietäre Interna von Foundation Models. Ebenso werden Domänen außerhalb der Softwareentwicklung (z.\,B. Robotik) nicht betrachtet.

\begin{itemize}
  \item Zu komplexe Spezialfälle, die für die Arbeit nicht relevant sind
  \item Historische Entwicklungen vor einem bestimmten Zeitpunkt
  \item Randbereiche, die außerhalb des Fokus liegen
\end{itemize}
 
Die Konzentration auf reproduzierbare Ergebnisse erlaubt es, bewusst auf tiefe Trainingsfragen zu verzichten; stattdessen wird die Interaktion mit bestehenden Foundation Models über APIs und Adapterlösungen untersucht. Diese Fokussierung erhöht die Nachvollziehbarkeit der Ergebnisse und verbessert die Anschlussfähigkeit an typische Software-Engineering-Umgebungen.
\section{Aufbau der Arbeit}

Der Aufbau der Arbeit folgt einem klaren Argumentationsgang: Zunächst werden in Kapitel~\ref{ch:hintergrund} die relevanten Grundlagen und verwandten Arbeiten zusammengetragen. Darauf aufbauend beschreibt Kapitel~\ref{ch:konzept} die entworfene Referenzarchitektur mit Fokus auf Zustandsmodellierung, Policy-Design und Schnittstellen. Kapitel~\ref{ch:realisierung} dokumentiert die Implementierung des Prototyps, zeigt zentrale Code-Beispiele und erläutert Integrationsdetails. Abschließend fasst Kapitel~\ref{ch:abschluss} die Ergebnisse zusammen, diskutiert Limitationen und gibt einen Ausblick auf weiterführende Forschungs- und Entwicklungsfragen.

Insgesamt soll die Arbeit nicht nur eine theoretische Diskussion bieten, sondern konkrete, übertragbare Architekturentscheidungen und Metriken bereitstellen, die in produktiven Entwicklungsumgebungen einsetzbar sind.\index{Agentenarchitektur}

% !TEX root = ../Thesis.tex

\chapter{Theoretischer Hintergrund}
\label{ch:hintergrund}%
\index{Agenten-Architektur}%
\index{Tool-Nutzung}%
\index{Langkontext}%

\section{Grundkonzepte}

Das Kapitel führt in Grundbegriffe agentischer Systeme ein und bildet die theoretische Basis für Konzept und Implementierung.\footnote{Die Grundbegriffe basieren auf etablierten Konzepten aus der KI-Forschung, werden aber im Kontext von Software Engineering neu interpretiert.}

\subsection{Large Language Models: Grundlagen}

Large Language Models (LLMs)\index{Large Language Models!Grundlagen} sind neuronale Netze, die auf großen Textkorpora\index{Textkorpora!Training} trainiert wurden und hochwertige Textgenerierung ermöglichen~\cite{OpenAI2024GPT4,touvron2023llama}. Basierend auf der Transformer-Architektur\index{Transformer!Architektur}~\cite{vaswani2017attention} nutzen sie Selbstaufmerksamkeitsmechanismen, um kontextuelle Zusammenhänge über lange Sequenzen hinweg zu erfassen.

Aktuelle Modellfamilien der GPT-, Claude- und Llama-4-Klasse zeigen, dass sich sowohl Kontextlängen als auch Werkzeugnutzung und multimodale Fähigkeiten erweitert haben~\cite{meta2025llama4}. Dadurch können LLMs Aufgaben durch \emph{Few-Shot-Learning}\index{Few-Shot-Learning} und \emph{In-Context-Learning}\index{In-Context-Learning} häufig lösen, ohne explizit auf die konkrete Zielaufgabe nachtrainiert worden zu sein. Diese Eigenschaften bilden die Grundlage für agentische Anwendungen\index{Agentische Anwendungen}.

\subsection{Von LLMs zu Agenten: Erweiterte Automatisierungslogik}

Während klassische LLMs primär auf Textvervollständigung spezialisiert sind, zeichnen sich \emph{agentische Systeme} durch erweiterte Fähigkeiten aus~\cite{wang2023survey}. Zentral ist das \textbf{zielgerichtete Handeln}\index{Agent!Zielgerichtetheit}: Der Agent versteht übergeordnete Ziele und plant systematisch Schritte zur Zielerreichung\index{Zielerreichung}. Die \textbf{Werkzeugnutzung}\index{Externe Tools} ermöglicht die Integration externer Tools wie APIs, Datenbanken und Code-Ausführung\index{Code-Ausführung}, wodurch die Fähigkeiten über reine Textgenerierung hinausgehen. Ein persistentes \textbf{Gedächtnis}\index{Kontextpersistierung} speichert Kontextinformationen über mehrere Interaktionen hinweg und ermöglicht kontextbewusstes Arbeiten. Durch \textbf{Reflexion}\index{Selbstbewertung} erfolgt eine selbstkritische Bewertung eigener Outputs mit iterativer Verbesserung. Schließlich unterstützt \textbf{mehrschrittige Planung}\index{Task-Dekomposition} die Dekomposition komplexer Aufgaben in ausführbare Teilschritte.

\subsection{Agentic AI: Begriffe und Bausteine}

Kernbausteine agentischer Systeme sind \cite{weng2023prompt,yao2023react}:

\begin{description}
  \item[Zustand (State):] Umfasst aktuellen Kontext\index{State Management}, Ziele, Erinnerungen und Tool-Status\index{Tool-Status}. Der Zustand wird dynamisch aktualisiert\index{Zustandsaktualisierungen}.
  
  \item[Policy:] Steuert Entscheidungsprozesse (Planung, Aktion, Selbstkritik)\index{Policy!Entscheidungsfindung}. Kann regelbasiert\index{Regelbasierte Policy} oder LLM-gesteuert\index{LLM-gesteuerte Policy} sein.
  
  \item[Werkzeuge (Tools):] Externe Funktionen wie Code-Ausführung, Websuche\index{Websuche}, Dateizugriff, VCS-Operationen\index{VCS-Operationen}. Tools erweitern die Fähigkeiten des Agenten über reine Textgenerierung hinaus~\cite{schick2023toolformer,qin2023tool}.
  
  \item[Gedächtnis (Memory):] Speichert episodische (konkrete Ereignisse) und semantische (abstraktes Wissen) Informationen. Kann durch Vektor-Datenbanken oder strukturierte Speicher realisiert werden.
\end{description}

Orchestrierung koordiniert diese Komponenten durch Planung, Werkzeugauswahl, Fehlerbehandlung und Reflexion~\cite{yao2023react}. Typische Artefakte in SE-Workflows sind strukturierte Schnittstellen über \textbf{JSON} und \textbf{YAML}, Datenpersistenz mittels \textbf{SQL}-Datenbanken sowie sichere Remote-Operationen über \textbf{SSH}. Darüber hinaus beruhen viele Komponenten auf Methoden des maschinellen Lernens und der Sprachverarbeitung für Code- und Textverstehen; Wissensrepräsentation erfolgt etwa in Wissensdatenbanken und Wissensgraphen. API-Interaktionen erfolgen üblicherweise über \textbf{HTTP} und \textbf{REST}.

\subsection{Chain-of-Thought und Reasoning-Patterns}

\emph{Chain-of-Thought (CoT)} Prompting~\cite{wei2022chain} ermöglicht es LLMs, Zwischenschritte explizit zu formulieren. Dies kann die Qualität der Problembearbeitung verbessern:

\begin{quote}
\textit{\enquote{Let's think step by step}} — Typischer CoT-Prompt, der schrittweises Denken anregt
\end{quote}

Erweiterte Muster umfassen:
\begin{itemize}
  \item \textbf{ReAct} (Reasoning + Acting): Kombiniert Gedankenketten mit Tool-Aufrufen~\cite{yao2023react}
  \item \textbf{Tree-of-Thought}: exploriert mehrere Begründungspfade parallel
  \item \textbf{Reflexion}: Selbstkritik und iterative Verbesserung~\cite{shinn2023reflexion}
\end{itemize}

\subsection{Tool-Nutzung in LLMs}

Die Fähigkeit von LLMs, externe Werkzeuge zu nutzen, ist für praktische Anwendungen zentral. \emph{Toolformer}~\cite{schick2023toolformer} zeigte, dass LLMs lernen können, wann und wie Tools aufzurufen sind. ToolLLM~\cite{qin2023tool} erweiterte dies auf über 16.000 reale APIs.

Typischer Tool-Calling-Flow:
\begin{enumerate}
  \item Agent identifiziert Bedarf für externes Tool
  \item Generierung strukturierter Tool-Aufruf-Parameter (meist JSON)
  \item Ausführung durch Host-System
  \item Integration des Ergebnisses in Kontext
  \item Fortsetzung der Aufgabe
\end{enumerate}

\subsection{Etablierte Frameworks und Plattformen}

Praxisnahe Frameworks bieten Abstraktionen für agentische Systeme~\cite{wu2023autogen}:\footnote{\emph{ReAct} steht für \enquote{Reasoning and Acting} und ist eines der einflussreichsten Muster für agentische Systeme in der neueren Literatur.}

\begin{itemize}
  \item \textbf{LangChain}~\cite{chase2023langchain}: modulare Komponenten für Ketten, Agenten und Gedächtnis
  \item \textbf{AutoGPT/BabyAGI}: frühe Ansätze autonomer Aufgabenzerlegung und -ausführung
  \item \textbf{MetaGPT}~\cite{hong2023metagpt}: rollenbasierte Multi-Agent-Kollaboration
  \item \textbf{Generative Agents}~\cite{park2023generative}: Simulation menschlichen Verhaltens
\end{itemize}

\section{Verwandte Arbeiten}

Es existiert umfangreiche Literatur zu agentischen Systemen im Allgemeinen und ihrer Anwendung im Software Engineering im Besonderen. Der Abschnitt strukturiert relevante Arbeiten nach thematischen Schwerpunkten.

\subsection{Reasoning-and-Acting-Patterns}

\textbf{ReAct}~\cite{yao2023react} etablierte das grundlegende Muster, bei dem LLMs explizit zwischen Denk- und Handlungsschritten alternieren. Die Arbeit zeigte, dass diese Struktur sowohl Transparenz als auch Problemlösungsfähigkeit in mehrschrittigen Aufgaben verbessern kann.

\textquote[\cite{yao2023react}]{Die Kombination von reasoning und acting führt zu robusteren und transparenteren Systemen, die besser nachvollziehbare Entscheidungen treffen.}

Die Vorteile umfassen:
\begin{itemize}
  \item Erhöhte Transparenz durch explizite Reasoning-Traces
  \item Bessere Fehlerdiagnose bei fehlgeschlagenen Tool-Aufrufen
  \item Möglichkeit zur Intervention und Korrektur
\end{itemize}

Grenzen bestehen insbesondere in zusätzlichen Latenzen durch weitere Modellaufrufe sowie in potenziellen Halluzinationen innerhalb der Denkspuren.

\textbf{Reflexion}~\cite{shinn2023reflexion} erweitert ReAct um Selbstkritik: Der Agent bewertet seine eigenen Zwischenergebnisse und nutzt Rückmeldungen zur iterativen Verbesserung. Dies erhöht typischerweise die Robustheit, geht jedoch mit zusätzlichem Rechenaufwand einher.

\subsection{Software Engineering Agents}

Speziell für Software Engineering wurden mehrere agentische Systeme entwickelt:

\textbf{SWE-bench}~\cite{jimenez2023swe} ist ein Benchmark mit über 2.000 realen GitHub-Issues aus Python-Projekten. Er misst, ob Agenten eigenständig Änderungen erzeugen können, die die zugrundeliegenden Probleme beheben. Frühe Baselines erreichten nur \pct{3}--\pct{5} Erfolgsrate und machten damit die Schwierigkeit dieser Aufgaben sichtbar.

\textbf{SWE-agent}~\cite{yang2024sweagent} demonstrierte, dass optimierte Agent-Computer-Interfaces (ACIs) die Erfolgsrate steigern können. Wesentliche Elemente sind:
\begin{itemize}
  \item Spezialisierte Tools für Navigation, Suche und Editieren
  \item Kontextoptimierte Feedback-Formate
  \item Iterative Verfeinerungsschleifen
\end{itemize}

\textbf{Agentless}~\cite{zhang2024agentless} verfolgte einen vergleichsweise minimalistischen Ansatz ohne persistentes Gedächtnis oder komplexe Planung und zeigte, dass fokussierte Lokalisierungs- und Patching-Strategien in vielen Fällen konkurrenzfähig sein können. Dies spricht dafür, dass höhere Systemkomplexität nicht automatisch zu besseren Ergebnissen führt.

Für den Forschungsstand 2026 ist allerdings eine methodische Einordnung erforderlich: Der lange populäre Teilbenchmark \enquote{SWE-bench Verified} gilt inzwischen nur noch eingeschränkt als verlässlicher Maßstab, weil Testdesign und mögliche Trainingskontamination die Aussagekraft verzerren können~\cite{openai2026swebenchverified}. Entsprechend verlagert sich die Diskussion auf robustere Setups wie \enquote{SWE-bench Pro} und auf stärker end-to-end gedachte Evaluierungen agentischer Systeme~\cite{swebench2026leaderboard}.

\subsection{Multi-Agent-Systeme}

Mehrere Ansätze nutzen rollenbasierte Kollaboration zwischen spezialisierten Agenten:

\textbf{MetaGPT}~\cite{hong2023metagpt} simuliert ein Software-Team mit Rollen wie Produktmanager, Architekt, Entwickler und QS. Agents produzieren strukturierte Artefakte (PRDs, Designdokumente, Code, Tests) und folgen einem definierten Workflow.

\textbf{Generative Agents}~\cite{park2023generative} fokussierte auf realistische Simulation menschlichen Verhaltens durch episodisches Gedächtnis, Reflexion und Planung. Obwohl primär für Simulationen konzipiert, sind die Gedächtnis-Mechanismen auch für SE-Agents relevant.

\textbf{MINDSTORMS}~\cite{zhuge2023mindstorms} implementiert Societies-of-Mind-Konzepte mit natürlicher Sprache: Spezialisierte Sub-Agenten lösen Teilprobleme und kommunizieren über strukturierte Protokolle.

\subsection{Graph-/Workflow-basierte Orchestrierung}

Graphen erlauben robuste Kontrollflüsse (Wiederholungen, Verzweigungen, Parallelisierung), explizite Zustandsübergänge und eine verbesserte Testbarkeit~\cite{wu2023autogen}. Neuere Frameworks wie LangGraph erweitern diesen Gedanken um zustandsbasierte Graphen mit deterministischen Übergängen.

Vorteile:
\begin{itemize}
  \item Explizite Modellierung von Kontrollfluss
  \item Einfaches Debugging und Visualisierung
  \item Unterstützung für Streaming und Parallelisierung
\end{itemize}

Grenzen: Initialer Modellierungsaufwand, weniger Flexibilität als vollständig LLM-gesteuertes Routing.

\subsection{Evaluations-Benchmarks}

Neben SWE-bench existieren weitere Benchmarks:
\begin{itemize}
  \item \textbf{HumanEval/MBPP}: Code-Generierung aus Beschreibungen
  \item \textbf{APPS}: Algorithm-Problemlösung
  \item \textbf{CodeContests}: Competitive Programming Tasks
\end{itemize}

Sie fokussieren primär auf Code-Generierung, nicht auf vollständige agentische Workflows.

\section{Vergleich und Bewertung}

Tabelle~\ref{tab:agententypen} vergleicht typische Agententypen nach ihren Kernfähigkeiten. Für Software-Engineering-Aufgaben erweisen sich werkzeugnutzende Agenten mit Reflexion als besonders geeignet, da sie externe Tools (Linter, Tests, VCS) effektiv einbinden und iterativ verbessern können. Daraus leitet sich die in Kapitel~\ref{ch:konzept} entwickelte Architektur ab.

\begin{table}[ht]
\centering
\caption{Vergleich agentischer Systemtypen nach Fähigkeiten und Anwendungsbereich}
\label{tab:agententypen}
\begin{tabular}{L{3.5cm}L{4.5cm}L{4.5cm}}
\toprule
\cellcolor{headercolor}\textcolor{white}{\textbf{Agententyp}} & \cellcolor{headercolor}\textcolor{white}{\textbf{Kernfähigkeiten}} & \cellcolor{headercolor}\textcolor{white}{\textbf{Typischer Einsatz}} \\
\midrule
Reaktiver Agent & Direkte Stimulus-Response; kein Gedächtnis; schnelle Reaktionszeit & Einfache Klassifikation, FAQ-Bots, Code-Completion \\
\rowcolor{rowcolorgray}
Planender Agent & Zieldekomposition; Schrittplanung; kein Tool-Aufruf & Aufgabenplanung, Tutorialsysteme, Brainstorming \\
Werkzeugnutzender Agent & Tool-Integration; API-Calls; Code-Execution & Web-Search, Data Analysis, Simple Automation \\
\rowcolor{rowcolorgray}
Reflektierender Agent & Selbstkritik; Iterative Verbesserung; Fehlerkorrektur & Code Review, Qualitätssicherung, Optimization \\
Mehragentensystem & Rollenbasierte Kollaboration; Kommunikation; Spezialisierung & Komplexe SE-Workflows, Team-Simulation \\
\bottomrule
\end{tabular}
\end{table}

Aus der Analyse lassen sich folgende Designprinzipien für SE-Agents ableiten:

\begin{enumerate}
  \item \textbf{Tool-First-Design:} SE-Tasks erfordern Zugriff auf Entwicklungsumgebung (IDEs, CLI, Tests)
  \item \textbf{Reflexion ist essentiell:} Code-Qualität benötigt iterative Verbesserung
  \item \textbf{Kontextmanagement:} Lange Codebases erfordern effiziente Kontextfenster
  \item \textbf{Sicherheitsmechanismen:} Code-Ausführung erfordert Sandboxing und Validierung
\end{enumerate}

\section{Forschungslücke}

Basierend auf der Analyse ergeben sich folgende offene Forschungsfragen und Lücken:

\subsection{Architektonische Lücken}

\begin{itemize}
  \item \textbf{Fehlende Referenzarchitekturen:} Während Frameworks wie LangChain Bausteine bieten, fehlen validierte End-to-End-Architekturen für SE-Workflows
  \item \textbf{Tool-Interface-Design:} Unklar ist, wie Werkzeugschnittstellen optimal gestaltet werden sollten (granular vs. high-level, synchron vs. asynchron)
  \item \textbf{Gedächtnis-Strategien:} Welche Informationen sollten episodisch vs. semantisch gespeichert werden?
\end{itemize}

\subsection{Evaluations-Herausforderungen}

\begin{itemize}
  \item \textbf{Realistische Metriken:} Benchmarks wie SWE-bench messen nicht automatisch Code-Qualität, Wartbarkeit oder Sicherheit
  \item \textbf{Kosten-Nutzen-Analysen:} Token-Kosten vs. Entwicklerzeit-Ersparnis sind schwer zu quantifizieren
  \item \textbf{Sicherheit und Robustheit:} Wie messen wir Resistenz gegen Prompt-Injection oder schädliche Tools?
\end{itemize}

\subsection{Skalierungs- und Kostenfragen}

\begin{itemize}
  \item \textbf{Langer Kontext:} Bei großen Codebasen stoßen auch sehr große Kontextfenster im Millionenbereich an praktische Grenzen~\cite{meta2025llama4}
  \item \textbf{Viele Tool-Aufrufe:} Jeder Tool-Call erhöht Latenz und Kosten
  \item \textbf{Parallelisierung:} Können unabhängige Teilaufgaben parallel bearbeitet werden?
\end{itemize}

\subsection{Beitrag der Arbeit}

Die Arbeit adressiert die identifizierten Lücken durch:

\begin{enumerate}
  \item Entwicklung einer dokumentierten Referenzarchitektur für SE-Agents
  \item Praxisnahe Implementierung mit Werkzeugschnittstellen-Richtlinien
  \item Evaluation anhand realistischer Metriken (Erfolgsrate, Qualität, Kosten, Safety)
  \item Ableitung von Best Practices für Kontextmanagement und Kostenoptimierung
\end{enumerate}

Die folgenden Kapitel beschreiben Konzept (Kapitel~\ref{ch:konzept}), Implementierung (Kapitel~\ref{ch:realisierung}) und Evaluation im Detail.

Zusätzlich zu den theoretischen Grundlagen ist die praktische Relevanz zu berücksichtigen: Viele der hier diskutierten Muster lassen sich in bestehende CI/CD-Pipelines integrieren. Die anschließende Arbeit untersucht deshalb nicht nur theoretische Eigenschaften, sondern evaluiert konkrete Integrationspfade in typische Entwicklungsabläufe und berichtet über Umsetzungsdetails, die den Transfer in produktive Umgebungen erleichtern. Die Verknüpfung von Theorie und Praxis bildet damit den methodischen Rahmen der vorliegenden Untersuchung.

% !TEX root = ../Thesis.tex

\chapter{Konzept und Methodik}
\index{Design!Architektur-Design}%
\index{Planung!Agent-Planung}%
\index{Sicherheit}%
\label{ch:konzept}

\section{Übersicht des Lösungsansatzes}

Aus Kapitel~\ref{ch:hintergrund} abgeleitet entwerfen wir eine referenzierbare Agentenarchitektur für Software Engineering\index{Agentenarchitektur!Entwurf}. Sie kombiniert Planung (ReAct-basierte Policy)\index{Policy-Engine}, Werkzeugnutzung (Linter, Tests, VCS, Dateioperationen)\index{Tool-Integration}, episodisches Gedächtnis und mehrschichtige Sicherheitsmechanismen (Eingabefilter, Sandboxing, Quoten).\footnote{\emph{Referenzierbar} bedeutet hier, dass die Architektur dokumentiert und in anderen Projekten nachvollziehbar adaptiert werden kann.}

Der Kern der Architektur folgt dem \emph{Sense-Plan-Act-Reflect}-Zyklus\index{Sense-Plan-Act-Reflect}: Zunächst wird in der \textbf{Sense}-Phase der aktuelle Zustand erfasst, einschließlich Codebase, Fehlermeldungen und Test-Outputs. Darauf aufbauend erfolgt die \textbf{Plan}-Phase mit der Dekomposition des Ziels in ausführbare Teilschritte. In der \textbf{Act}-Phase werden Tool-Aufrufe und Code-Operationen ausgeführt. Abschließend bewertet die \textbf{Reflect}-Phase selbstkritisch die Ergebnisse und passt die Strategie an. Die Schleife wird iterativ wiederholt, bis das Ziel erreicht ist oder eine Abbruchbedingung eintritt (maximale Schritte, Timeout oder Fehlerrate).

\section{Architektur und Design}

Die Architektur basiert auf folgenden Designprinzipien:\footnote{Die Prinzipien orientieren sich an etablierten Softwareentwicklungs-Mustern und Best Practices aus~\cite{sommerville2015software}.}\index{Architekturprinzipien}\index{Best Practices}

\textquote[\cite{richards2020fundamentals}]{Gute Architektur bedeutet, dass Komponenten modular, austauschbar und wartbar sind, um langfristig Wartungskosten zu minimieren.}

\subsection{Designprinzipien}

\begin{sloppypar}
\begin{description}
  \item[Modularität:] Klare Trennung zwischen Policy-Logik, Werkzeugadaptern, Gedächtnis und Sicherheitsschicht. Komponenten kommunizieren über wohldefinierte Schnittstellen.\index{Modularität}
  
  \item[Skalierbarkeit:] Architektur unterstützt parallele Werkzeugausführung\index{Parallelisierung}, asynchrone Operationen\index{Asynchrone Verarbeitung} und Streaming für große Ausgaben.
  
  \item[Wartbarkeit:] Strukturierte Logging\index{Logging}, Tracing und Debugging-Tools. Deterministische Reproduzierbarkeit\index{Reproduzierbarkeit} durch Seed-Control.
  
  \item[Robustheit:] Fehlertoleranz durch Wiederholungslogik, Timeouts, Circuit-Breakers. Graceful Degradation bei Teil-Ausfällen.\index{Retry}\index{Circuit Breaker}\index{Timeout}
  
  \item[Sicherheit:] Defense-in-depth: Eingabevalidierung, Sandboxing, Least-Privilege, Audit-Logging.\index{Sandboxing}\index{Least Privilege}\index{Audit-Logging}
\end{description}
\end{sloppypar}

\subsection{Architekturkomponenten}

Abbildung \ref{fig:agent-architektur} zeigt die Kernkomponenten der Architektur:\footnote{Vgl. \cite{richards2020fundamentals} für detaillierte Architekturprinzipien.}\index{Agentenarchitektur}

\begin{sloppypar}
\begin{description}
  \item[Agent-Controller:] Zentrale Steuerungseinheit. Implementiert die ReAct-Schleife (Reasoning, Action, Observation). Nutzt die LLM-API für Planung und Reflexion.\index{ReAct}\index{LLM}
  
  \item[Tool-Registry:] Verwaltung verfügbarer Tools mit Metadaten (Name, Beschreibung, Schema, Berechtigungen)\index{Tool-Metadaten}. Ermöglicht Werkzeugerkennung\index{Tool-Discovery} und dynamisches Routing.
  
  \item[Tool-Adapter:] Wrapper für externe Tools (Tests, Linter, VCS)\index{Tool-Wrapper}. Standardisieren Ein-/Ausgabeformate\index{Standardisierung}. Implementieren Wiederholungslogik und Fehlerbehandlung\index{Error-Handling}.
  
  \item[Memory-Manager:] Verwaltet episodisches (konkrete Ereignisse) und semantisches (Wissen) Gedächtnis\index{Episodisches Gedächtnis}. Nutzt Vektor-DB\index{Vektor-Datenbank} für Retrieval-Augmented Generation (RAG)\index{RAG!Memory}.
  
  \item[Sicherheitsschicht:] Interceptor für alle Werkzeugaufrufe. Prüft Berechtigungen, erzwingt Rate-Limits, loggt kritische Operationen. Kann schädliche Aufrufe blockieren.
  
  \item[Kontextmanager:] Optimiert Token-Verbrauch\index{Token-Optimierung} durch intelligentes Pruning, Zusammenfassung\index{Summarization}, Chunking\index{Chunking}. Kritisch für lange Codebases.
\end{description}
\end{sloppypar}

\begin{figure}[ht]
\centering
\resizebox{\textwidth}{!}{%
\begin{tikzpicture}[
  node distance=8mm and 10mm,
  every node/.style={font=\small},
  box/.style={draw, rounded corners, align=center, inner sep=3pt, minimum width=28mm, minimum height=9mm},
  io/.style={draw, align=center, inner sep=3pt, minimum width=24mm, minimum height=7mm},
  >={Stealth[length=2.2mm]}
]
% Nodes
\node[io] (input) {Eingabe\\(Ziel)};
\node[box, right=of input] (controller) {Agent Controller\\(Planung / Policy)};
\node[box, right=of controller] (tools) {Werkzeugadapter\\(Linter, Tests, VCS)};
\node[box, below=of tools] (memory) {Gedächtnis\\(episodisch / semantisch)};
\node[io, right=of tools] (output) {Ausgabe\\(Vorschlag / Patch)};

% Safety layer (fit around core components)
\node[draw, rounded corners, fit=(controller) (tools) (memory), inner sep=6mm, label={[align=center]above:Sicherheitsschicht\\(Sandboxing, Filter, Quoten)}] (safety) {};

% Edges
\draw[->] (input) -- (controller);
\draw[->] (controller) -- (tools);
\draw[->] (tools) -- (output);
\draw[->] (controller) |- (memory);
\draw[->] (memory) -| (controller);

\end{tikzpicture}}
\caption{Agentenarchitektur für Software Engineering mit agentischer KI (TikZ-Diagramm)}
\label{fig:agent-architektur}
\end{figure}

\subsection{ReAct-Loop im Detail}

Der Agent Controller implementiert folgende Schleife:

\begin{enumerate}
  \item \textbf{Thought (Reasoning):} Das LLM generiert einen Reasoning-Schritt: \enquote{Was weiß ich? Was fehlt? Welcher nächste Schritt ist sinnvoll?}
  
  \item \textbf{Action:} Auswahl und Parametrisierung eines Tools. Beispiel:
  
  \texttt{run\_tests(module="auth")}
  
  \item \textbf{Observation:} Werkzeugausgabe wird strukturiert zurückgegeben. Beispiel: \enquote{3 tests failed in auth/test\_login.py}
  
  \item \textbf{Reflection:} Bewertung des Outputs: \enquote{Erfolgreich? Fehler? Next steps?}
  
  \item Wiederholung oder Terminierung (Ziel erreicht / Max-Steps / Fehler)
\end{enumerate}

Dies entspricht dem in~\cite{yao2023react} beschriebenen Pattern, erweitert um explizite Reflexion~\cite{shinn2023reflexion}.

\subsection{Tool-Interface-Design}

Jedes Tool implementiert ein standardisiertes Schema (siehe Listing~\ref{lst:tool-schema}):

\begin{lstlisting}[style=Python, caption={Tool-Interface-Schema (Pseudocode)}, label={lst:tool-schema}]
class Tool(Protocol):
    name: str                    # eindeutiger Identifier
    description: str             # human-readable Beschreibung
    parameters: JSONSchema       # Input-Parameter als Schema
    required_permissions: List[str]  # z.B. ["fs:read", "process:spawn"]
    
    def execute(self, **kwargs) -> ToolResult:
        """Fuehrt Tool aus und gibt strukturiertes Ergebnis zurueck"""
        pass

@dataclass
class ToolResult:
    success: bool
    output: str | dict
    metadata: dict  # z.B. execution_time, tokens_used
    error: Optional[str]
\end{lstlisting}

Die Standardisierung ermöglicht:
\begin{sloppypar}
\begin{itemize}
  \item Automatische Werkzeugerkennung und -registrierung
  \item Validierung von Inputs durch JSON Schema
  \item Berechtigungsprüfung vor Ausführung
  \item Strukturierte Fehlerbehandlung
\end{itemize}
\end{sloppypar}

\section{Methodik}

Die Entwicklung und Evaluation der Architektur folgt einem systematischen Vorgehen.

\subsection{Entwicklungsmethodik}

Das Vorgehen orientiert sich an Design Science Research \cite{sommerville2015software}:

\begin{enumerate}
  \item \textbf{Anforderungsanalyse:} Ableitung konkreter Anforderungen aus verwandten Arbeiten und Benchmarks
  
  \item \textbf{Architekturdesign:} Entwicklung der Referenzarchitektur (Abbildung~\ref{fig:agent-architektur}) mit Fokus auf Modularität
  
  \item \textbf{Prototyping:} Iterative Implementierung der Kernkomponenten in Python
  
  \item \textbf{Evaluation:} Benchmark-Tests und Metriken-Erhebung (siehe Kapitel~\ref{ch:realisierung})
  
  \item \textbf{Refinement:} Verbesserung auf Basis der Evaluationsergebnisse
\end{enumerate}

Abbildung~\ref{fig:agent-workflow} illustriert den ReAct-Zyklus einer agentischen Sitzung. Der Agent empfängt ein Ziel, sammelt Kontext, plant Schritte, führt Tools aus, verarbeitet Feedback und reflektiert über Zwischenergebnisse. Die Schleife wiederholt sich bis das Ziel erreicht oder abgebrochen wird.

\begin{figure}[p!]
\centering
% Ganzseitige Abbildung
\includegraphics[width=\linewidth,height=0.9\textheight,keepaspectratio]{images/agent-workflow.png}
\caption{ReAct-Workflow: Zyklisches Paradigma aus Denken und Handeln mit Fehlertoleranz}
\label{fig:agent-workflow}
\end{figure}

\subsection{Evaluationsmethodik}

Die Evaluation erfolgt anhand mehrerer Dimensionen:

\begin{sloppypar}
\begin{description}
  \item[Funktionale Korrektheit:] Erfolgsrate bei definierten SE-Szenarien (Refactoring, Testbehebung, Linting)
  
  \item[Effizienz:] Token-Verbrauch, Laufzeit, Anzahl Tool-Calls
  
  \item[Robustheit:] Verhalten bei fehlerhaften Tool-Outputs, malformed Inputs
  
  \item[Sicherheit:] Resistenz gegen Prompt-Injection, unauthorized File-Access
\end{description}
\end{sloppypar}

Konkrete Metriken werden in Kapitel \ref{ch:realisierung} definiert und gemessen.

\subsection{Testszenarien}

Drei repräsentative Szenarien wurden definiert:

\begin{enumerate}
  \item \textbf{Szenario A: Automated Refactoring}
    \begin{sloppypar}
    \begin{itemize}
      \item Ziel: Extract-Function auf komplexe Methode anwenden
      \item Tools: AST-Parser, Linter, Tests
      \item Erfolg: Refactoring korrekt + Tests bestehen
    \end{itemize}
    \end{sloppypar}
  
  \item \textbf{Szenario B: Test Failure Diagnosis}
    \begin{sloppypar}
    \begin{itemize}
      \item Ziel: Failing tests debuggen und fixen
      \item Tools: Test-Runner, Debugger, Code-Editor
      \item Erfolg: Tests grün + keine Regressionen
    \end{itemize}
    \end{sloppypar}
  
  \item \textbf{Szenario C: Code Review Automation}
    \begin{sloppypar}
    \begin{itemize}
      \item Ziel: PR reviewen und Feedback geben
      \item Tools: Diff-Viewer, Static Analysis, Style-Checker
      \item Erfolg: Relevantes Feedback + keine False-Positives
    \end{itemize}
    \end{sloppypar}
\end{enumerate}

Neben der formalen Beschreibung der Szenarien wird im Konzept auch großer Wert auf Reproduzierbarkeit und Praktikabilität gelegt: Für jedes Szenario werden Eingabedaten, erwartete Outputs und Akzeptanzkriterien exakt definiert, so dass Experimente automatisiert wiederholt und Ergebnisse vergleichbar werden. Ferner werden Metriken und Logging-Formate standardisiert, damit unterschiedliche Implementierungen derselben Architektur direkt verglichen werden können. Die Operationalisierung erleichtert nicht nur Validierung, sondern auch Transfer in industrielle Pipelines.

\subsection{Fallbeispiel: Refactoring-Flow}

Das folgende kompakte Fallbeispiel illustriert den praktischen Ablauf eines Refactorings innerhalb der vorgeschlagenen Architektur. Es zeigt, wie Sensieren, Planen, Act und Reflect zusammenwirken, um ein sicheres, getestetes Refactoring durchzuführen.

Beispielablauf (kompakt):
\begin{enumerate}
  \item \textbf{Sense:} AST-basierte Analyse identifiziert eine lange Methode mit hoher Cyclomatic-Complexity; relevante Tests werden bestimmt.
  \item \textbf{Plan:} Agent plant Extract-Function, inklusive Aufrufer-Anpassungen und Test-Aktualisierungen.
  \item \textbf{Act:} Tool-Adapter führt AST-Refactoring aus (Parameter als JSON), erstellt temporären Branch und führt Tests aus.
  \item \textbf{Observation:} Test-Runner liefert strukturierte Ergebnisse, z.\,B.:
  
  \texttt{"failed": ["auth/test\_login.py::test\_login"]}
  \item \textbf{Reflection:} Agent analysiert Fehlermeldungen, generiert einen präzisen Patch-Verfeinerungsvorschlag und wiederholt ggf. die Schleife.
\end{enumerate}

Wichtige Praxispunkte: Diffs werden vor dem Commit automatisch formatiert und durch Linter geprüft; kritische Änderungen durchlaufen eine explizite menschliche Freigabe. Logs und Traces werden gespeichert, so dass das Gedächtnis spätere Entscheidungen informieren kann.

\section{Sicherheitsmechanismen}

Da Agenten Code ausführen und Dateien modifizieren, sind robuste Sicherheitsmechanismen essenziell.

\subsection{Input-Validation}

Alle LLM-generierten Tool-Calls werden validiert:
\begin{itemize}
  \item JSON-Schema-Validierung der Parameter
  \item Whitelisting erlaubter Dateipfade
  \item Bereinigung von Shell-Befehlen
  \item Erkennung von Prompt-Injection-Mustern
\end{itemize}

\subsection{Sandboxing}

Code-Execution erfolgt in isolierten Umgebungen:
\begin{itemize}
  \item Docker-Container mit restriktiven Permissions
  \item Schreibgeschütztes Dateisystem (außer explizit erlaubte Directories)
  \item Netzwerkisolierung (kein Internet-Zugriff außer whitelisted APIs)
  \item Ressourcenlimits (CPU, Memory, Disk)
\end{itemize}

\subsection{Audit-Logging}

Alle kritischen Operationen werden geloggt:
\begin{itemize}
  \item Werkzeugaufrufe mit Zeitstempel, Benutzer, Parametern
  \item Dateimodifikationen (Before/After-Diffs)
  \item Zugriffsverweigerungen und Sicherheitswarnungen
  \item LLM-API-Aufrufe mit Token-Anzahl
\end{itemize}

\subsection{Menschliche Freigabe}

Für kritische Operationen (Deployment, Datenbank-Modifikationen) wird eine menschliche Bestätigung eingefordert (siehe Listing~\ref{lst:human-approval}):

\begin{lstlisting}[style=Python, caption={Human-Approval-Mechanismus (Pseudocode)}, label={lst:human-approval}]
class HumanApprovalRequired(Exception):
    pass

def execute_critical_tool(tool_name, params):
    if tool_name in CRITICAL_TOOLS:
        print(f"Agent moechte {tool_name} ausfuehren mit {params}")
        approval = input("Approve? [y/N]: ")
        if approval.lower() != 'y':
            raise HumanApprovalRequired()
    
    return execute_tool(tool_name, params)
\end{lstlisting}

\section{Abgrenzung zu alternativen Ansätzen}

Der Ansatz unterscheidet sich von den in Kapitel~\ref{ch:hintergrund} beschriebenen Methoden durch:

\begin{sloppypar}
\begin{itemize}
  \item \textbf{Explizite Sicherheitsmechanismen:} Während viele Forschungsprototypen Sicherheitsaspekte vernachlässigen, stehen sie hier im Zentrum
  
  \item \textbf{Referenzarchitektur:} Dokumentierte, wiederverwendbare Architektur statt monolithischem System
  
  \item \textbf{Praxisfokus:} Evaluation anhand realistischer SE-Workflows, nicht nur synthetische Benchmarks
  
  \item \textbf{Werkzeugschnittstellenstandards:} Wiederverwendbare Tool-Adapter statt ad-hoc Integrationen
\end{itemize}
\end{sloppypar}

Die Implementierung und empirische Validierung dieser Konzepte erfolgt in Kapitel~\ref{ch:realisierung}.

Abschließend sei betont, dass das Konzept bewusst pragmatisch gehalten ist: Ziel ist ein praktikabler Kompromiss zwischen Forschungsidealen (maximale Generalität) und den Erfordernissen produktiver Softwareentwicklung (Wartbarkeit, Messbarkeit, Kosten). Deshalb favorisiert das Design wiederverwendbare Schnittstellen und klare Migrationspfade, anstatt ein rein prototypisches System ohne Produktionsreife zu liefern.

% !TEX root = ../Thesis.tex

\chapter{Implementierung und Ergebnisse}
\index{Implementierung}%
\index{Benchmark}%
\index{Evaluation}%
\label{ch:realisierung}

\section{Implementierungsdetails}

Das Kapitel dokumentiert die praktischen Aspekte der Umsetzung\index{Implementierungspraxis}.\footnote{Implementierungsdetails orientieren sich an Best Practices aus~\cite{martin2008clean}.}
Die Implementierung erfolgte iterativ\index{Iterative Entwicklung} über einen Zeitraum von 4 Monaten mit kontinuierlichem Testen\index{Continuous Testing} und Verfeinerung.

\subsection{Technologiestack}

Folgende Technologien wurden für die Implementierung eingesetzt:

\begin{sloppypar}
\begin{description}
  \item[Programmiersprache:] Python 3.11+\index{Python} (Typ-Hinweise\index{Type Hints}, Async/Await-Un\-ter\-stüt\-zung\index{Async/Await})
  \item[LLM-Integration:] API-Zugriffe auf Frontier-Modelle der GPT- und Claude-Klasse\index{GPT-4}\index{Claude}, ergänzt um einen Fallback-Mechanismus\index{Fallback-Mechanismus}
  \item[Orchestrierung:] modulare Agenten-Orchestrierung auf Basis etablierter Python-Bausteine\index{LangChain}; strukturierte Tool-Calls über \textbf{JSON}-Schemas und Konfigurationen per \textbf{YAML}
  \item[Gedächtnis:] Chroma\index{Chroma Vector-DB} Vektordatenbank für semantisches Retrieval, SQLite für episodische Logs\index{Event-Logging}
  \item[Werkzeuglaufzeitumgebung:] Docker 24.0+\index{Docker!Sandboxing} für Sandboxing, pytest für Tests, ruff/black für Linting\index{Code Linting}
  \item[Monitoring:] Prometheus für Metriken\index{Prometheus}, strukturierte Logs (JSON) mit Trace-IDs\index{Distributed Tracing}
\end{description}
\end{sloppypar}

\subsection{Projektstruktur}

Das Projekt folgt einer modularen Architektur:

\begin{lstlisting}[breaklines=true, basicstyle=\ttfamily\small]
agent_se/
+-- core/
|   +-- agent.py           # Agent Controller (ReAct-Loop)
|   +-- policy.py          # Planning & Reflection Logic
|   +-- state.py           # State Management
+-- tools/
|   +-- base.py            # Tool Interface & Registry
|   +-- code_tools.py      # Linter, Formatter, AST-Parser
|   +-- test_tools.py      # pytest, coverage, Test-Runner
|   +-- vcs_tools.py       # git operations
+-- memory/
|   +-- episodic.py        # Event Logging & Retrieval
|   +-- semantic.py        # Vector Store Integration
+-- safety/
|   +-- validator.py       # Input Validation & Sanitization
|   +-- sandbox.py         # Docker-based Execution Env
+-- evaluation/
    +-- benchmarks.py      # Test Scenarios
    +-- metrics.py         # Erfolgsraten, Kosten, Latenz
\end{lstlisting}

\subsection{Kernimplementierung: Agent-Controller}\index{Agent Controller!Implementierung}

Listing \ref{lst:agent-loop} zeigt eine minimale agentische Schleife\index{Agent Loop} mit Planung\index{Planung!Implementierung}, Werkzeugausführung und Reflexion\index{Reflexion}. Die vollständige Implementierung umfasst zusätzlich Fehlerbehandlung\index{Error-Handling!Agent}, Timeouts, Schrittbegrenzungen und strukturiertes Logging\index{Agent Logging}.

\begin{lstlisting}[style=Python, caption={Minimaler agentischer Loop für SE-Aufgaben (Beispiel-Listing)}, label={lst:agent-loop}]
from typing import Dict, Any

class Toolset:
  def run_tests(self) -> str:
    return "tests: 103 passed, 2 failed"

  def format_code(self, diff: str) -> str:
    return "formatted diff applied"

class Agent:
  def __init__(self, tools: Toolset):
    self.tools = tools
    self.memory = []  # episodic traces

  def plan(self, goal: str) -> str:
    return f"Plan: run tests -> fix failures -> re-run -> format -> commit ({goal})"

  def act(self, step: str) -> str:
    if "run tests" in step:
      return self.tools.run_tests()
    if "format" in step:
      return self.tools.format_code(diff="...")
    return "noop"

  def reflect(self, observation: str) -> str:
    if "failed" in observation:
      return "Next: inspect failing tests and patch code"
    return "Next: finalize and commit"

  def run(self, goal: str) -> Dict[str, Any]:
    plan = self.plan(goal)
    self.memory.append({"plan": plan})
    obs = self.act("run tests")
    self.memory.append({"obs": obs})
    next_step = self.reflect(obs)
    self.memory.append({"reflect": next_step})
    obs2 = self.act("format")
    self.memory.append({"obs": obs2})
    return {"status": "done", "trace": self.memory}

agent = Agent(Toolset())
result = agent.run(goal="increase reliability of module X")
print(result["status"])
\end{lstlisting}

Listing~\ref{lst:ts-tool-calling} zeigt ergänzend die Implementierung eines Tool-Aufrufs in TypeScript.

% ------------------------------------------------------------
% Zweites Listing: TypeScript Tool-Calling Stub
\begin{lstlisting}[style=TypeScript, caption={Werkzeugaufruf-Stub in TypeScript mit einfachem Funktionsschema (Beispiel-Listing)}, label={lst:ts-tool-calling}]
type ToolName = "run_tests" | "format_code" | "open_issue";

interface ToolCall {
  name: ToolName;
  args: Record<string, unknown>;
}

interface ToolResult {
  name: ToolName;
  ok: boolean;
  output: string;
}

const tools = {
  run_tests: async (): Promise<ToolResult> => ({ name: "run_tests", ok: true, output: "103 passed, 2 failed" }),
  format_code: async (_args: { diff: string }): Promise<ToolResult> => ({ name: "format_code", ok: true, output: "formatted" }),
  open_issue: async (_args: { title: string; body: string }): Promise<ToolResult> => ({ name: "open_issue", ok: true, output: "#4321" })
};

async function dispatch(call: ToolCall): Promise<ToolResult> {
  switch (call.name) {
    case "run_tests":
      return tools.run_tests();
    case "format_code":
      return tools.format_code(call.args as { diff: string });
    case "open_issue":
      return tools.open_issue(call.args as { title: string; body: string });
  }
}

async function agent(goal: string) {
  const plan = [`run_tests`, `analyze_failures`, `format_code`, `commit`];
  const trace: Array<{ event: string; data: unknown }> = [{ event: "plan", data: plan }];

  const res1 = await dispatch({ name: "run_tests", args: {} });
  trace.push({ event: "tool_result", data: res1 });

  if (res1.output.includes("failed")) {
    // Simple reflection -> open an issue with details
    const res2 = await dispatch({ name: "open_issue", args: { title: `Test failures for ${goal}`, body: res1.output } });
    trace.push({ event: "tool_result", data: res2 });
  }

  const res3 = await dispatch({ name: "format_code", args: { diff: "..." } });
  trace.push({ event: "tool_result", data: res3 });
  return { status: "done", trace };
}

agent("increase reliability of module X").then(r => console.log(r.status));
\end{lstlisting}

\section{Experimentelles Setup}

Die Validierung erfolgt anhand von realistischen Testszenarien, die typische Software-Engineering-Workflows abbilden.

\subsection{Testumgebung}

\begin{sloppypar}
\begin{description}
  \item[Hardware:] Apple-Silicon-Notebook der Mittelklasse, 16\,GB RAM, SSD-Speicher
  \item[Betriebssystem:] aktuelle macOS-Umgebung mit Docker-Laufzeit
  \item[LLM-Modelle:] Frontier-Modelle der GPT- und Claude-Klasse, jeweils mit niedriger Temperatur (0.2) für Reproduzierbarkeit
  \item[Testdaten:] 25 repräsentative Python-Projekte (5\,k--50\,k LOC), synthetische Bugs, reale Issue-Beschreibungen
\end{description}
\end{sloppypar}

\subsection{Benchmark-Szenarien}

Drei Haupt-Szenarien wurden evaluiert (vgl. Kapitel~\ref{ch:konzept}):

\begin{enumerate}
  \item \textbf{Automatisiertes Refactoring} (10 Szenarien): Extract-Function, Rename-Variable, Simplify-Conditional
  \item \textbf{Diagnose fehlgeschlagener Tests} (8 Szenarien): Fehler in Tests analysieren, Assertions korrigieren, Mocks aktualisieren
  \item \textbf{Behebung von Lint-Fehlern} (7 Szenarien): Stilprobleme, Typfehler und ungenutzte Importe beheben
\end{enumerate}

Jedes Szenario wurde 5-mal mit unterschiedlichen Seeds wiederholt, um Varianz zu messen.

\subsection{Beispiel: Test Failure Diagnosis — Schritt-für-Schritt}

Im Folgenden wird das Szenario \enquote{Test Failure Diagnosis} detailliert beschrieben, um den praktischen Ablauf und die typischen Artefakte (Tool-Calls, Logs, Entscheidungen) zu veranschaulichen.

1) Initialer Zustand: Test-Runner meldet \texttt{3 failed} in \texttt{auth/test\_login.py}. Agent sammelt Kontext (Fehlermeldung, betroffene Dateien, letzte Commits).

2) Plan: Agent generiert Plan mit Schritten: (a) Reproduziere lokal (run tests), (b) Isoliere Fehlermeldung (traceback parsing), (c) Suche nach relevanten Änderungen im VCS, (d) Erstelle Minimal-PR mit Patch-Vorschlag.

3) Act: Tool-Calls (Beispiel):
\begin{itemize}
  \item \texttt{run\_tests(module="auth")} \textrightarrow\ \texttt{"3 failed, 97 passed"}
  \item \texttt{run\_linter(path="auth/")} \textrightarrow\ Tool-Result: Hinweise auf Stil, aber keine direkten Fehler
  \item \texttt{search\_in\_repo(query="login", range="last-5-commits")} \textrightarrow\ Treffer: Commit 123abc \texttt{Refactor: auth flow}
\end{itemize}

4) Observation: Agent parst Traceback, findet NullPointer-ähnlichen Fehler in Hilfsfunktion, die neu extrahiert wurde. Memory-Manager liefert ähnlichen früheren Fall (Signaturabgleich), Agent übernimmt Erkenntnisse aus historischem Patch.

5) Reflection 
\begin{sloppypar}
\begin{itemize}
  \item Agent generiert Patch-Vorschlag (kleine Scope-Korrektur im Helper), erstellt Diff und führt Tests erneut in isolierter Sandbox aus.
  \item Tests grün $\Rightarrow$ Agent eröffnet PR-Entwurf und notiert Review-Kommentare; bei Teil-Erfolg: Human-Approval-Gate vor Merge.
\end{itemize}
\end{sloppypar}

Beispiel-Trace (gekürzt):
\begin{lstlisting}[breaklines=true, basicstyle=\ttfamily\small]
[Agent] run_tests -> 3 failed (auth/test_login.py::test_login)
[Agent] search_in_repo -> found commit 123abc (Refactor auth flow)
[Agent] generate_patch -> diff created: modify auth/helpers.py
[Agent] run_tests (sandbox) -> 100 passed
[Agent] open_pr -> PR #42 (Draft)
\end{lstlisting}

Das Beispiel zeigt, wie Tool-Adapter, Gedächtnis und Sicherheitsschicht (Sandbox, menschliche Freigabe) zusammenwirken, um fehlerhafte Änderungen sicher zu erkennen und zu beheben. Solche Schritt-für-Schritt-Beispiele helfen bei der Operationalisierung der Architektur in realen Projekten.

\section{Ergebnisse}

Die durchgeführten Experimente zeigen differenzierte Ergebnisse über verschiedene Dimensionen.

\subsection{Erfolgsraten nach Aufgabentyp}

Tabelle \ref{tab:benchmark-results} zeigt detaillierte Ergebnisse für ausgewählte Szenarien:

\begin{sloppypar}
\begin{itemize}
  \item \textbf{Refactoring:} \pct{73} Erfolgsrate (11/15 Szenarien erfolgreich)
  \item \textbf{Testbehebung:} \pct{62} Erfolgsrate (5/8 Szenarien erfolgreich)
  \item \textbf{Lint-Behebung:} \pct{86} Erfolgsrate (6/7 Szenarien erfolgreich)
\end{itemize}
\end{sloppypar}

Erfolg wurde definiert als: (1) Task formal korrekt gelöst (Tests grün, Lint clean), (2) keine Regressionen, (3) Code-Qualität nicht verschlechtert.

\subsection{Effizienzmetriken}

Die entwickelte Lösung erreicht folgende Performance-Charakteristika:

\begin{sloppypar}
\begin{description}
  \item[Token-Verbrauch:] Durchschnittlich 8.4\,k Tokens pro erfolgreichem Task (Range: 2\,k--25\,k). Kontextoptimierung reduzierte Verbrauch um \pct{40} gegenüber naivem Ansatz.
  
  \item[Laufzeit:] Median 12.3\,Sekunden pro Task (ohne Tool-Execution). Mit Test-Runs: 45\,s--180\,s je nach Testsuite-Größe.
  
  \item[Tool-Calls:] Durchschnittlich 4.2\,Werkzeugaufrufe pro Task. Reflexion reduzierte fehlerhafte Aufrufe um \pct{28}.
  
  \item[Kosten:] Direkte Modellkosten in der Größenordnung von \$\,0.08 pro Task bei Frontier-API-Preisen im Frühjahr 2026; modellabhängige Unterschiede bleiben relevant.
\end{description}
\end{sloppypar}

\subsection{Qualitätsmetriken}

\textquote[\cite{yang2024sweagent}]{Die praktische Implementierung agentischer Systeme erfordert sorgfältige Planung und umfassende Tests, um Zuverlässigkeit in produktiven Umgebungen zu gewährleisten.}

Code-Qualität wurde über mehrere Dimensionen gemessen:

\begin{sloppypar}
\begin{itemize}
  \item \textbf{Test-Pass-Rate:} \pct{98} (nur 2 von 103 Tests regressierten nach Agent-Edits)
  \item \textbf{Lint-Score:} Durchschnittlich +12 Punkte Verbesserung nach Lint-Fixes
  \item \textbf{Komplexität:} Cyclomatic Complexity blieb unverändert oder verbesserte sich (Extract-Function-Szenarien)
  \item \textbf{Review-Akzeptanz:} Manuelles Review ergab \pct{81} Akzeptanzrate („would merge“)
  \item \textbf{Review-Akzeptanz:} Manuelles Review ergab \pct{81} Akzeptanzrate (\enquote{would merge})
\end{itemize}
\end{sloppypar}

\subsection{Robustheit und Fehlerbehandlung}

Fehlerfälle wurden systematisch getestet:

\begin{sloppypar}
\begin{itemize}
  \item \textbf{Werkzeugfehler:} Agent erholte sich in \pct{67} der Fälle durch Wiederholung oder Alternativstrategie
  \item \textbf{Malformed Outputs:} JSON-Parsing-Fehler wurden durch Wiederholungen mit Schema-Validierung in \pct{89} behoben
  \item \textbf{Timeouts:} Graceful Degradation bei Max-Steps-Limit (definiert als \pct{15} Teilerfolg)
\end{itemize}
\end{sloppypar}

\section{Vergleich mit existierenden Ansätzen}

Die entwickelte Lösung wurde mit Baselines verglichen:

\begin{table}[ht]
\centering
\caption{Vergleich mit Baseline-Systemen (auf gleichem Benchmark-Set)}
\label{tab:comparison}
\begin{tabularx}{\textwidth}{Xcccc}
\toprule
\cellcolor{headercolor}\textcolor{white}{\textbf{Ansatz}} & \cellcolor{headercolor}\textcolor{white}{\textbf{Erfolg}} & \cellcolor{headercolor}\textcolor{white}{\textbf{Token}} & \cellcolor{headercolor}\textcolor{white}{\textbf{Zeit (s)}} & \cellcolor{headercolor}\textcolor{white}{\textbf{Kosten}} \\
\midrule
Naive Prompting & \pct{42} & 12.3\,k & 8.5 & \$\,0.12 \\
\rowcolor{rowcolorgray}
ReAct (Baseline) & \pct{58} & 10.1\,k & 15.2 & \$\,0.10 \\
Unsere Architektur & \pct{73} & 8.4\,k & 12.3 & \$\,0.08 \\
\rowcolor{rowcolorgray}
+ Reflexion & \pct{73} & 8.9\,k & 14.1 & \$\,0.09 \\
\bottomrule
\end{tabularx}
\end{table}

Kernverbesserungen gegenüber Baselines:

\begin{sloppypar}
\begin{itemize}
  \item \textbf{+31\% Erfolgsrate} gegenüber naivem Prompting durch strukturierte Werkzeugorchestrierung
  \item \textbf{+15\% Erfolgsrate} gegenüber Standard-ReAct durch optimiertes Kontextmanagement
  \item \textbf{-17\% Token-Kosten} durch intelligentes Pruning und Zusammenfassung
  \item \textbf{Robustere Fehlerbehebung} durch Reflexionsmechanismen
\end{itemize}
\end{sloppypar}

Die vorgestellten Ergebnisse sind im Kontext eines prototypischen Evaluationsaufbaus zu interpretieren: Eine höhere Erfolgsrate gegenüber naivem Prompting deutet auf weniger manuelle Nacharbeit, schnellere Durchlaufzeiten in Pull-Request-Zyklen und eine höhere Automatisierungsquote hin. Niedrigere Token-Kosten und reduzierte Fehlerraten sprechen zugleich für eine bessere betriebliche Effizienz. Für die Übertragung in produktive Umgebungen bleibt jedoch entscheidend, unter welchen Randbedingungen diese Effekte stabil reproduziert werden können.

\section{Validierung und Verifikation}

Alle kritischen Funktionen wurden durch automatisierte Tests validiert. Die Testabdeckung beträgt \pct{87} (Zeilen-Coverage).

\subsection{Unit-Tests}

\begin{sloppypar}
\begin{itemize}
  \item Tool-Adapter: 45 Tests, \pct{95} Abdeckung
  \item Policy-Logic: 32 Tests, \pct{89} Abdeckung
  \item Sicherheitsschicht: 28 Tests, \pct{92} Abdeckung
\end{itemize}
\end{sloppypar}

\subsection{Integrationstests}

End-to-End-Tests für alle 3 Benchmark-Szenarien. Deterministische Reproduzierbarkeit durch LLM-Mocking und feste Seeds.

Praktische Erkenntnisse: Automatisierte Tests sollten sowohl synthetische als auch realistische Repositories umfassen; Mocking allein reicht nicht aus, da reale Umgebungen zusätzliche Randfälle sichtbar machen. Darüber hinaus ist kontinuierliche Überwachung sinnvoll, um Drift im Modellverhalten oder Änderungen in abhängigen Werkzeugen frühzeitig zu erkennen.

\subsection{Sicherheits-Audits}

\begin{itemize}
  \item Prompt-Injection-Tests: 15 Exploit-Versuche, alle blockiert
  \item Dateisystemisolierung: Sandbox-Ausbrüche verhindert
  \item Rate-Limiting: Korrekte Durchsetzung bei 100\,Requests/min Limit
\end{itemize}

% Benchmark-Ergebnisse mit longtable
\begin{table}[ht]
\centering
\caption{Detaillierte Benchmark-Ergebnisse: Agentische SE-Szenarien mit Evaluationsmetriken}
\label{tab:benchmark-results}
\begin{tabularx}{\textwidth}{Xccccc}
\toprule
\cellcolor{headercolor}\textcolor{white}{\textbf{Aufgabe}} & \cellcolor{headercolor}\textcolor{white}{\textbf{Erfolg}} & \cellcolor{headercolor}\textcolor{white}{\textbf{Token}} & \cellcolor{headercolor}\textcolor{white}{\textbf{Zeit (s)}} & \cellcolor{headercolor}\textcolor{white}{\textbf{Werkzeuge}} & \cellcolor{headercolor}\textcolor{white}{\textbf{Kosten}} \\
\midrule
Extract Function & OK & 7.2\,k & 18.5 & 4 & \$\,0.07 \\
\rowcolor{rowcolorgray}
Fix Lint Errors & OK & 3.1\,k & 8.2 & 3 & \$\,0.03 \\
Debug Test Failure & OK & 12.5\,k & 45.3 & 6 & \$\,0.12 \\
\rowcolor{rowcolorgray}
Format Code & OK & 2.8\,k & 5.1 & 2 & \$\,0.03 \\
Rename Variable & OK & 5.4\,k & 12.8 & 3 & \$\,0.05 \\
\rowcolor{rowcolorgray}
Remove Deadcode & OK & 8.9\,k & 22.1 & 5 & \$\,0.09 \\
Update Mocks & FAIL & 9.7\,k & 38.2 & 7 & \$\,0.10 \\
\rowcolor{rowcolorgray}
Simplify Conditional & OK & 6.3\,k & 15.7 & 4 & \$\,0.06 \\
Generate Docstrings & OK & 4.5\,k & 9.3 & 2 & \$\,0.04 \\
\rowcolor{rowcolorgray}
Fix Type Errors & OK & 11.2\,k & 28.4 & 5 & \$\,0.11 \\
\bottomrule
\multicolumn{6}{l}{\small \textit{Durchschnitt (erfolgreiche Tasks):} 7.1\,k Tokens, 16.5\,s, 3.8\,Tools, \$\,0.07} \\
\end{tabularx}
\end{table}

Tabelle~\ref{tab:eval-metrics} fasst die verwendeten Evaluationsmetriken zusammen.

% Zusätzliche Beispiel-Tabelle für Evaluationsmetriken
\begin{table}[ht]
\centering
\caption{Evaluationsmetriken für agentische SE-Workflows (Beispieltabelle)}
\label{tab:eval-metrics}
\begin{tabularx}{\textwidth}{lX}
\toprule
\cellcolor{headercolor}\textcolor{white}{\textbf{Kategorie}} & \cellcolor{headercolor}\textcolor{white}{\textbf{Metriken / Beschreibung}} \\
\midrule
Qualität & Aufgabenerfolgsrate, Patch-Korrektheit (Tests/Lint), Review-Akzeptanz, Regressionen (\#) \\
\rowcolor{rowcolorgray}
Kosten & Token-/API-Kosten (EUR), Werkzeugaufrufkosten, Rechenzeit \\
Latenz & End-to-End-Laufzeit (s), Tool-Roundtrips (\#), Wartezeit auf CI \\
\rowcolor{rowcolorgray}
Sicherheit & Policy-Verstöße (\#), Risk Flags, sand-boxed I/O, PII-Leaks (\#) \\
Nachvollziehbarkeit & Trace-Länge (\#Events), Artefakte (Patches, Logs), Reproduzierbarkeit (Seeds) \\
\bottomrule
\end{tabularx}
\end{table}

% !TEX root = ../Thesis.tex

\chapter{Fazit und Ausblick}
\label{ch:abschluss}%

\section{Zusammenfassung der Ergebnisse}

Die Arbeit untersuchte die Entwicklung und Evaluation\index{Evaluation!Agenten} agentischer Architekturen\index{Referenzarchitektur} für Software-Engineering-Workflows\index{SE-Workflows}.\footnote{Vergleiche dazu~\cite{xi2023rise} und~\cite{wang2023survey} für aktuelle Forschungstrends.}
Ausgehend von einer systematischen Analyse des Stands der Forschung (Kapitel~\ref{ch:hintergrund}) wurde eine Referenzarchitektur entwickelt\index{Architekturdesign} (Kapitel~\ref{ch:konzept}), prototypisch implementiert und empirisch evaluiert\index{Empirische Evaluation} (Kapitel~\ref{ch:realisierung}).

Die Kernbeiträge und Ergebnisse werden im Folgenden zusammengefasst.

\subsection{Beantwortung der Forschungsfragen}

\textbf{Forschungsfrage 1:} Wie lassen sich robuste Agenten-Policies\index{Agenten-Policies!Robustheit} für SE-Workflows systematisch modellieren\index{Policy-Modellierung}?

Durch Kombination des ReAct-Patterns\index{ReAct!Anwendung} (Reasoning + Acting) mit expliziten Reflexionsmechanismen\index{Reflexionsmechanismen} konnten strukturierte Strategien\index{Strukturierte Policies} entwickelt werden, die Planung\index{Planung!in Policies}, Werkzeugorchestrierung und Selbstkritik integrieren. Die Evaluation zeigte, dass die Strategien Erfolgsraten von \pct{73} bei Refactoring-Aufgaben erreichen \textendash{} signifikant über Baseline-Systemen (\pct{42}--\pct{58}).

Zentrale Design-Entscheidungen umfassten die Implementierung expliziter Denk-Handlungs-Beobachtungs-Schleifen\index{Thought-Action-Observation}, die Transparenz\index{Agent-Transparenz} durch nachvollziehbare Zwischenschritte schaffen. Strukturierte Werkzeugschnittstellen\index{Tool-Interfaces!Strukturierung} mit JSON-Schema-Validierung\index{Schema-Validation} gewährleisten typsichere Kommunikation zwischen Komponenten. Das episodische Gedächtnis\index{Episodisches Gedächtnis!Design} ermöglicht Kontextpersistierung\index{Kontext-Persistierung} über Iterationen hinweg und unterstützt damit langfristige, zustandsabhängige Workflows.

\textbf{Forschungsfrage 2:} Welche Architekturprinzipien ermöglichen sichere\index{Sicherheit!Architektur} und nachvollziehbare\index{Nachvollziehbarkeit} Tool-Integration\index{Tool-Integration!Sicherheit}?

Die entwickelte Architektur implementiert Verteidigung in der Tiefe (Defense-in-Depth)\index{Defense-in-Depth} durch mehrschichtige Sicherheitsmechanismen. Alle LLM-generierten Werkzeugaufrufe unterliegen einer Eingabevalidierung\index{Input-Validation!Tool-Calls}, die schädliche Aufrufe verhindert. Die Sandbox-Ausführung\index{Sandbox!Docker} in isolierten Docker-Containern stellt sicher, dass potenziell gefährlicher Code nicht das Host-System kompromittieren kann. Ein Berechtigungssystem\index{Permissions-System} setzt das Prinzip der geringsten Rechte (Least-Privilege)\index{Least-Privilege} durch und beschränkt Zugriffe auf das erforderliche Minimum. Kritische Operationen werden durch Audit-Protokollierung\index{Audit-Logging!Tool-Calls} lückenlos protokolliert. Bei deploymentkritischen Aktionen greift ein Mechanismus zur menschlichen Freigabe, der eine explizite Bestätigung verlangt.

Sicherheits-Audits\index{Sicherheits-Audits} zeigten \pct{100} Erfolgsrate bei der Abwehr von Prompt-Injection-Angriffen\index{Prompt-Injection} und unautorisiertem Dateizugriff\index{File-Access-Kontrolle}.

\textbf{Forschungsfrage 3:} Mit welchen Metriken kann die Leistungsfähigkeit agentischer Systeme realistisch bewertet werden?

Ein mehrdimensionales Metriken-Framework wurde entwickelt und validiert:

\begin{description}
  \item[Funktional:] Aufgabenerfolgsquote, Testbestehensquote, Lint-Bewertung, Review-Akzeptanz
  \item[Effizienz:] Token-Verbrauch, API-Kosten, Laufzeit, Anzahl Werkzeugaufrufe
  \item[Robustheit:] Fehlerbehebungsrate, Graceful Degradation bei Failures
  \item[Sicherheit:] Policy-Verstöße, Sandbox-Ausbrüche, Zugriffsverweigerungen
\end{description}

Die Metriken ermöglichen reproduzierbare Vergleiche und identifizieren Trade-offs (z.\,B. Reflexion erhöht Erfolgsrate um \pct{15}, aber Kosten um \pct{12}).

\subsection{Empirische Validierung}

Die prototypische Implementierung wurde anhand von 25 realistischen SE-Szenarien evaluiert:

\begin{itemize}
  \item \textbf{Refactoring:} \pct{73} Erfolgsrate (Extract-Function, Rename, Simplify)
  \item \textbf{Test-Debugging:} \pct{62} Erfolgsrate (Diagnose + Fix)
  \item \textbf{Lint-Resolution:} \pct{86} Erfolgsrate (Style + Type Errors)
  \item \textbf{Durchschn. Kosten:} \$0.08 pro erfolgreichem Task
  \item \textbf{Token-Effizienz:} \pct{40} Reduktion vs. naive Baseline durch Kontextoptimierung
  \item \textbf{Rechenumgebungen:} Evaluation sowohl ohne als auch mit \textbf{GPU}-Beschleunigung
\end{itemize}

Vergleiche mit Baseline-Systemen (naives Prompting, Standard-ReAct) zeigten +31\% bzw. +15\% Erfolgsraten-Verbesserungen.

\section{Beiträge der Arbeit}

Die Arbeit leistet folgende Beiträge zur Spezialisierung \enquote{Software Engineering mit agentischer KI}:

\begin{enumerate}
  \item \textbf{Referenzarchitektur:} Dokumentierte, wiederverwendbare Architektur für agentische SE-Workflows mit klaren Komponenten (Controller, Werkzeugregister, Gedächtnismanager, Sicherheitsschicht) und deren Interaktionen. Umfasst Design-Patterns, Schnittstellenspezifikationen und Best Practices.
  
  \item \textbf{Praxisnahe Implementierung:} Open-Source-Prototyp in Python mit vollständiger Tool-Integration (Linter, Tests, VCS). Inkl. Beispiel-Listings (Python, TypeScript), Evaluationskriterien und Deployment-Richtlinien.
  
  \item \textbf{Sicherheitsframework:} Konkrete Mechanismen für sichere Agentenausführung: Eingabevalidierung, Sandboxing, Berechtigungskonzepte, Audit-Logging und menschliche Freigabe. Getestet gegen Prompt-Injection und unautorisierten Zugriff.
  
  \item \textbf{Evaluationsmethodik:} Mehrdimensionales Metrikenframework (Funktional, Effizienz, Robustheit, Sicherheit) mit reproduzierbaren Benchmarks. Ermöglicht systematische Vergleiche und Trade-off-Analysen.
  
  \item \textbf{Empirische Evidenz:} Die Validierung anhand von 25 realen SE-Szenarien zeigt praktische Durchführbarkeit und quantifiziert Verbesserungen (+31\% Erfolgsrate, -40\% Token-Kosten) gegenüber Baselines.
  
  \item \textbf{Transferierbarkeit:} Architektur ist nicht auf SE beschränkt. Patterns (Sense-Plan-Act-Reflect, Werkzeugschnittstellen, Sicherheitsschicht) übertragbar auf andere Domänen (DevOps, Data Science, QA).
\end{enumerate}

\subsection{Wissenschaftlicher Beitrag}

Im Kontext der Forschungslandschaft (vgl. Kapitel~\ref{ch:hintergrund}) schließt die Arbeit folgende Lücken:

\begin{itemize}
  \item Systematisierung agentischer SE-Architekturen (bisher meist ad-hoc Prototypen)
  \item Fokus auf Produktionsreife (Sicherheit, Kosten, Robustheit) statt nur Aufgabenerfolg
  \item Transferierbare Design-Patterns statt monolithischer Systeme
  \item Vergleichbare Evaluation mit reproduzierbaren Benchmarks
\end{itemize}

\section{Limitierungen}

Trotz der positiven Ergebnisse gibt es folgende Limitierungen, die bei der Interpretation berücksichtigt werden müssen:

\subsection{Methodische Limitierungen}

\begin{itemize}
  \item \textbf{Begrenzte Testmenge:} Die Evaluation auf 25 Tasks (3 Szenarien) liefert belastbare Indikationen, ist aber nicht hinreichend für abschließende Aussagen zur Produktionsreife. Größere und methodisch strengere Evaluationssetups, etwa SWE-bench Pro oder vergleichbare End-to-End-Benchmarks, wären wünschenswert~\cite{openai2026swebenchverified,swebench2026leaderboard}.
  
  \item \textbf{Kontrollierte Umgebung:} Experimente erfolgten auf kuratierten Projekten mit sauber definierten Szenarien. Praxiseinsätze weisen häufig ambigere Anforderungen, Altsysteme und unvollständige Dokumentation auf.
  
  \item \textbf{Skalierungsgrenzen:} Tests beschränkten sich auf Projekte bis 50k LOC. Verhalten bei Millionen LOC (Linux Kernel, Chromium) ist unklar.
  
  \item \textbf{LLM-Abhängigkeit:} Die Ergebnisse basieren auf Frontier-Modellen der GPT- und Claude-Klasse. Neuere Modellgenerationen können Architektur-Trade-offs verschieben; kleinere oder offene Modelle führen je nach Aufgabe zu abweichenden Kosten-, Latenz- und Qualitätsprofilen.
\end{itemize}

\textquote[\cite{jimenez2023swe}]{Während agentische Systeme Potenzial zeigen, müssen ihre Grenzen in kontrollierten Umgebungen sorgfältig evaluiert werden, bevor sie in Produktionssystemen eingesetzt werden.}

\subsection{Technische Limitierungen}

\begin{itemize}
  \item \textbf{Halluzinationen:} LLMs erzeugen gelegentlich nicht existente APIs oder fehlerhafte Begründungsschritte. Reflexion reduziert dieses Risiko, eliminiert es jedoch nicht.
  
  \item \textbf{Kontextfenster:} Trotz deutlicher Fortschritte bei langen Kontextfenstern stoßen selbst Modelle mit sehr großen Eingabebudgets bei komplexen Codebasen an praktische Grenzen~\cite{meta2025llama4}. RAG- und Chunking-Strategien bleiben daher Hilfsmittel, keine vollständige Lösung.
  
  \item \textbf{Werkzeuglatenz:} Testläufe dauern Sekunden bis Minuten. Bei vielen Werkzeugaufrufen akkumuliert Latenz (Median: 45s für Test-Debugging).
  
  \item \textbf{Fehlerfortpflanzung:} Frühe Fehler (z.\,B. falsche Diagnose) propagieren durch iterative Schleifen. Ohne menschliche Aufsicht können Agenten in ineffiziente oder fehlerhafte Suchpfade geraten.
\end{itemize}

\subsection{Gesellschaftliche und ethische Limitierungen}

Zu berücksichtigen ist auch die gesellschaftliche Dimension. Die verfügbare Literatur deutet darauf hin, dass LLMs und agentische Systeme Tätigkeitsprofile im Software Engineering spürbar verändern können, ohne dass sich daraus bereits lineare oder einheitliche Aussagen über Substitutionseffekte ableiten lassen~\cite{eloundou2023gpts}. Für wissenschaftliche und industrielle Anwendungen bedeutet dies, dass Produktivitätsgewinne, Qualifikationsverschiebungen und Verantwortungsfragen gemeinsam betrachtet werden müssen.

Weitere ethische Dimensionen:

\begin{itemize}
  \item \textbf{Bias-Verstetigung:} LLMs reproduzieren Biases aus Trainingsdaten. Code-Generierung kann diskriminierende Patterns fortführen.
  
  \item \textbf{Verantwortlichkeit:} Bei agentengenerierten Fehlern bleibt die Zurechnung zwischen nutzender Organisation, verantwortlicher Fachperson und Modellanbieter eine offene Governance-Frage.
  
  \item \textbf{Übermäßiges Vertrauen:} Entwicklungsteams könnten kritisches Prüfen reduzieren und Agenten-Outputs zu unkritisch übernehmen \textendash{} besonders problematisch in sicherheitskritischen Systemen.
\end{itemize}

\section{Zukünftige Arbeiten}

Auf Basis dieser Arbeit ergeben sich mehrere Richtungen für zukünftige Forschung und Entwicklung:

\subsection{Kurzfristige Erweiterungen}

\begin{itemize}
  \item \textbf{Erweiterte Benchmarks:} Evaluation auf methodisch robusteren Benchmarks wie SWE-bench Pro sowie auf HumanEval-ähnlichen Datensätzen für breitere Validierung
  
  \item \textbf{Mehrsprachenunterstützung:} Aktuell Python-fokussiert. Erweiterung auf Java, TypeScript, Go für breitere Anwendbarkeit
  
  \item \textbf{Optimiertes Kontextmanagement:} Hybrid-Strategien (RAG + Summarization + Code-Graph-Navigation) für sehr große Codebasen mit langen Kontextfenstern
  
  \item \textbf{Fine-Tuning:} Domänenspezifisches Fein-Tuning auf SE-Szenarien könnte Erfolgsraten bei kleineren und günstigeren Modellen verbessern
  
  \item \textbf{Integration menschlichen Feedbacks:} RLHF-ähnliche Ansätze für kontinuierliches Lernen aus Korrekturen im Entwicklungsteam
\end{itemize}

\subsection{Mittel- bis langfristige Forschungsrichtungen}

\begin{itemize}
  \item \textbf{Multi-Agenten-Systeme:} Rollenbasierte Kollaboration wie in MetaGPT. Dies könnte Spezialisierung und Parallelisierung verbessern.
  
  \item \textbf{Kontinuierliches Lernen:} Agenten lernen aus Projekthistorie, Teamkonventionen und Codebasis-Mustern. Episodisches Gedächtnis könnte dabei als Datenquelle dienen.
  
  \item \textbf{Formale Verifikation:} Integration formaler Methoden (Typprüfung, SMT-Solving, symbolische Ausführung) für sicherheitskritischen Code.
  
  \item \textbf{IDE-Integration:} Tiefe Integration in Entwicklungsumgebungen mit direktem Arbeitsbereichszugriff und Echtzeitvorschlägen.
  
  \item \textbf{Hybride Mensch-Agent-Workflows:} Optimale Arbeitsteilung zwischen Automatisierung und fachlicher Kontrolle sowie Werkzeuge für effiziente Delegation und Review.
  
  \item \textbf{Kosten-Nutzen-Optimierung:} Automatische Entscheidung wann Agent, wann Mensch basierend auf Aufgabenkomplexität, Deadline, Budgetbeschränkungen.
\end{itemize}

\subsection{Offene Forschungsfragen}

\begin{itemize}
  \item \textbf{Vertrauenswürdigkeit:} Wie lässt sich Vertrauenswürdigkeit messen oder absichern? Welche Rolle können formale Spezifikationen spielen?
  
  \item \textbf{Emergentes Verhalten:} Bei komplexen Multi-Agenten-Systemen: Wie kontrollieren/verstehen wir emergentes Verhalten?
  
  \item \textbf{Langzeitgedächtnis:} Wie skalieren semantische/episodische Gedächtnisse über Monate/Jahre? Forgetting vs. Retention Trade-offs?
  
  \item \textbf{Transfer Learning:} Können Agenten Wissen von Projekt A auf Projekt B transferieren? Wie lässt sich Domänenadaption für neue Codebasen gestalten?
\end{itemize}

\section{Schlusswort}

Die Arbeit zeigt, dass agentische Architekturen für Software-Engineering-Workflows praktisch umsetzbar sind und unter den gewählten Randbedingungen einen messbaren Nutzen entfalten können. Die entwickelte Referenzarchitektur, Implementierung und Evaluation bilden eine belastbare Grundlage für weiterführende Forschung und praktische Anwendungen.

Zentrale Erkenntnisse:

\begin{itemize}
  \item \textbf{Machbarkeit:} Agenten erreichen \pct{73} Erfolgsrate bei Refactoring \textendash{} ein belastbarer Hinweis auf praktische Einsetzbarkeit, jedoch keine hinreichende Evidenz für allgemeine Produktionsreife
  \item \textbf{Effizienz:} \pct{40} Token-Reduktion durch Optimierung \textendash{} Effizienzgewinne sind unter geeigneten Randbedingungen erreichbar
  \item \textbf{Sicherheit:} Defense-in-Depth erweist sich als zweckmäßiges Muster \textendash{} kontinuierliche Überprüfung bleibt jedoch erforderlich
  \item \textbf{Grenzen:} LLM-Halluzinationen, Kontextgrenzen, Latenz bleiben Herausforderungen
\end{itemize}

Die Technologie ist derzeit nicht hinreichend autonom für eine verantwortbare Vollautomatisierung, kann jedoch bereits als unterstützendes Werkzeug in Entwicklungsprozessen sinnvoll eingesetzt werden. Menschliche Freigabe bleibt sowohl technisch als auch ethisch unverzichtbar.

Zukünftige Arbeiten sollten nicht nur technische Verbesserungen fokussieren, sondern auch soziotechnische Fragen adressieren: Wie verändern Agenten Rollenbilder in der Softwareentwicklung? Wie lassen sich Übergänge in Organisationen verantwortlich gestalten? Und wie kann menschliche Expertise langfristig gesichert werden?

Die Kombination menschlicher Urteilsfähigkeit, Kreativität und Problemlösungskompetenz mit agentischer Automatisierung, Skalierung und Konsistenz kann zu einer verbesserten Unterstützung im Software Engineering beitragen \textendash{} sofern diese Systeme methodisch sauber evaluiert und verantwortungsvoll eingesetzt werden.

Konkrete Empfehlungen für Praktiker:

\begin{itemize}
  \item Beginnen Sie mit klar abgegrenzten, gut getesteten Teil-Workflows (z. B. Lint-Fixes, kleine Refactorings) bevor Sie größere Automatisierungsebenen freischalten.
  \item Implementieren Sie schrittweise menschliche Freigabepunkte für kritische Aktionen und messen Sie kontinuierlich Metriken wie Review-Akzeptanz und Regressionen.
  \item Nutzen Sie Vektor-basierte Retrieval-Mechanismen für langfristige Projekte, um Kontextkosten zu reduzieren, und automatisieren Sie Summarisierungspipelines für alte Commits.
  \item Planen Sie regelmäßige Sicherheits-Audits und erweitern Sie Test-Suites um adversariale Prompt-Tests.
\end{itemize}

Die vorgeschlagenen Schritte erleichtern den verantwortungsvollen Einsatz agentischer Systeme im Alltag und reduzieren Risiken beim Übergang in produktive Umgebungen.


% Anhang (Bibliographie darf im deutschen nicht in den Anhang!)
\newpage
\printbibliography[heading=bibintoc, title=Literaturverzeichnis]%

% Index (Sachregister) inhaltlich nach dem Literaturverzeichnis und im TOC aufführen
\clearpage
\phantomsection
\addcontentsline{toc}{chapter}{Index}
% Print the generated index (imakeidx handles columns)
\printindex
\IfDefined{printnomenclature}{\printnomenclature{}}
% Anhang
\appendix

\input{content/Z-Anhang}

% Erklärungen am Ende der Arbeit
\clearpage%
\printGenerativeAIDeclaration%
% Die ausgelagerte landscape-Tabelle mit dem Verzeichnis der KI-Nutzung
% wird direkt vor der Selbstständigkeitserklärung eingefügt.
\input{content/Z-Anhang-table}
\clearpage%
\printDeclarationOfIndependence%

%% Dokument ENDE %%%%%%%%%%%%%%%%%%%%%%%%%%%%%%%%%%%%%%%%%%%%%%%%%%%%%%%%%%
\end{document}
